% ==================================================
% Product Rule and Calculus Rules
% Discrete Calculus
% ==================================================

\section{Calculus Rules}
\label{sec:calculus-rules}

With the difference operator and the polynomial basis in place, we can
now develop the full suite of differentiation-like rules for discrete
calculus. The linearity of \(\Delta\) was established in the Difference
Operators section (Proposition~\ref{prop:delta-linearity}). Here we add
the \emph{product rule} — the discrete counterpart of the Leibniz rule
\((fg)' = f'g + fg'\).

\section*{Product Rule}

\begin{tcolorbox}[colback=thmbox, colframe=thmborder, arc=2pt,
  left=6pt, right=6pt, top=4pt, bottom=4pt,
  title={\small\textbf{Theorem (Discrete Product Rule)}},
  fonttitle=\small\bfseries]
For functions \(f, g : \mathbb{Z} \to \mathbb{R}\),
\[
\Delta(fg)(n)
=
f(n)\,\Delta g(n) + g(n+1)\,\Delta f(n).
\]
\end{tcolorbox}
\label{thm:product-rule}

\begin{remark*}[Fully quantified form]
\(\forall f,g : \mathbb{Z}\to\mathbb{R},\; \forall n \in \mathbb{Z}:\;
(f(n+1)g(n+1)) - (f(n)g(n))
= f(n)(g(n+1)-g(n)) + g(n+1)(f(n+1)-f(n))\).
\end{remark*}

\begin{proof}
By definition of \(\Delta\),
\[
\Delta(fg)(n) = f(n+1)g(n+1) - f(n)g(n).
\]
Add and subtract \(f(n)g(n+1)\):
\[
= f(n+1)g(n+1) - f(n)g(n+1) + f(n)g(n+1) - f(n)g(n).
\]
Grouping the first two and last two terms:
\[
= g(n+1)\bigl(f(n+1)-f(n)\bigr) + f(n)\bigl(g(n+1)-g(n)\bigr)
= g(n+1)\,\Delta f(n) + f(n)\,\Delta g(n). \qedhere
\]
\end{proof}

\begin{remark*}[Comparison with the continuous product rule]
In continuous calculus, \((fg)' = f'g + fg'\) is symmetric in \(f\)
and \(g\). In discrete calculus the rule is slightly asymmetric: one
factor is evaluated at \(n\) and the other at \(n+1\). This asymmetry
is unavoidable and is a characteristic feature of the discrete setting.
\end{remark*}

\begin{remark*}[Consequence — summation by parts]
The product rule leads directly to the summation by parts formula:
summing both sides of \(\Delta(fg) = f\,\Delta g + g(E)\,\Delta f\)
over \(k = a\) to \(b-1\) and applying the Fundamental Theorem gives
the formula in the next section.
\end{remark*}
